% TurnAhead 论文的中文阅读版主文件（对应 ACL Rolling Review 投稿稿）。
% 编译：latexmk -xelatex main.tex
% !TEX program = xelatex
\documentclass[11pt]{article}

% 允许参考文献中的长网址在连字符处换行。
\PassOptionsToPackage{hyphens}{url}
% ARR 投稿使用 review；后续可改为 final 或 preprint。
\usepackage[review]{acl}

\usepackage{times}
\usepackage{latexsym}
\usepackage{microtype}
\usepackage{graphicx}

\usepackage{amsmath}
\usepackage{amssymb}
\usepackage{amsthm}
% 中文阅读版的共享排版设置；使用 XeLaTeX 编译。
\usepackage{fontspec}
\usepackage{xeCJK}
\setmainfont{NimbusRoman-Regular.otf}[Path=/usr/share/fonts/opentype/urw-base35/,BoldFont=NimbusRoman-Bold.otf,ItalicFont=NimbusRoman-Italic.otf,BoldItalicFont=NimbusRoman-BoldItalic.otf]
\setsansfont{NimbusSans-Regular.otf}[Path=/usr/share/fonts/opentype/urw-base35/,BoldFont=NimbusSans-Bold.otf,ItalicFont=NimbusSans-Italic.otf,BoldItalicFont=NimbusSans-BoldItalic.otf]
\setmonofont{DejaVu Sans Mono}[Scale=MatchLowercase]
\setCJKmainfont{Noto Serif CJK SC}[BoldFont={Noto Serif CJK SC Bold},ItalicFont={Noto Serif CJK SC}]
\setCJKsansfont{Noto Sans CJK SC}[BoldFont={Noto Sans CJK SC Bold}]
\setCJKmonofont{Noto Sans Mono CJK SC}[BoldFont={Noto Sans Mono CJK SC Bold},Scale=0.9908]
\renewcommand{\abstractname}{摘要}
\renewcommand{\figurename}{图}
\renewcommand{\tablename}{表}
\renewcommand{\appendixname}{附录}
\renewcommand{\refname}{参考文献}
\providecommand{\proofname}{证明}
\renewcommand{\proofname}{证明}
% ACL 样式中的匿名作者栏和参考文献标题采用固定英文文本，在此本地化。
\usepackage{etoolbox}
\makeatletter
\patchcmd{\outauthor}{Anonymous ACL submission}{ACL 匿名投稿}{}{}
\patchcmd{\thebibliography}{References\@mkboth{References}{References}}{参考文献\@mkboth{参考文献}{参考文献}}{}{}
\makeatother
\emergencystretch=1em

\usepackage{booktabs}
\usepackage{multirow}
\usepackage{enumitem}
\usepackage{xspace}
\usepackage{xcolor}
\usepackage{algorithm}
\usepackage{algpseudocode}
\usepackage{tikz}
\usetikzlibrary{arrows.meta,positioning,calc,fit,backgrounds,decorations.pathreplacing}
% cleveref 必须在 acl.sty 加载的 hyperref 之后加载。
\usepackage[capitalize,noabbrev]{cleveref}

% !TEX root = main.tex
% 名称。策略名称为暂定名称，仅在此处统一修改。
\newcommand{\method}{TurnAhead\xspace}
\newcommand{\spork}{SPORK\xspace}
\newcommand{\stratLaunch}{锚定前瞻\xspace}
\newcommand{\stratCommit}{精确输入提交\xspace}
\newcommand{\stratRecover}{排空后恢复\xspace}

% 智能体循环的记号（与第 2 节共同定稿）。
\newcommand{\hist}{H}                % 已提交历史
\newcommand{\ctx}[1]{X_{#1}}         % 第 #1 轮的模型输入
\newcommand{\render}{F}              % 历史策略／提示词渲染器
\newcommand{\rsn}[1]{r_{#1}}         % 推理文本
\newcommand{\vis}[1]{v_{#1}}         % 可见的非推理文本
\newcommand{\act}[1]{a_{#1}}         % 动作（工具调用）
\newcommand{\obs}[1]{o_{#1}}         % 观测
\newcommand{\env}[1]{S_{#1}}         % 环境状态
\newcommand{\anchor}{\mathcal{A}}    % 锚点
\newcommand{\slack}{\sigma}          % 锚点就绪后的剩余时间
\newcommand{\contract}{I}            % 提交契约

% 定理类环境。
\theoremstyle{plain}
\newtheorem{proposition}{命题}
\newtheorem{lemma}{引理}
\newtheorem{corollary}{推论}
\theoremstyle{definition}
\newtheorem{definition}{定义}
\newtheorem{assumption}{假设}
\theoremstyle{remark}
\newtheorem{remark}{注}
\crefname{section}{节}{节}
\Crefname{section}{节}{节}
\crefformat{section}{#2第~#1~节#3}
\Crefformat{section}{#2第~#1~节#3}
\crefname{subsection}{节}{节}
\Crefname{subsection}{节}{节}
\crefformat{subsection}{#2第~#1~节#3}
\Crefformat{subsection}{#2第~#1~节#3}
\crefname{appendix}{附录}{附录}
\Crefname{appendix}{附录}{附录}
\crefname{equation}{式}{式}
\Crefname{equation}{式}{式}
\crefname{figure}{图}{图}
\Crefname{figure}{图}{图}
\crefname{table}{表}{表}
\Crefname{table}{表}{表}
\crefname{algorithm}{算法}{算法}
\Crefname{algorithm}{算法}{算法}
\crefname{proposition}{命题}{命题}
\Crefname{proposition}{命题}{命题}
\crefname{lemma}{引理}{引理}
\Crefname{lemma}{引理}{引理}
\crefname{corollary}{推论}{推论}
\Crefname{corollary}{推论}{推论}
\crefname{definition}{定义}{定义}
\Crefname{definition}{定义}{定义}
\crefname{assumption}{假设}{假设}
\Crefname{assumption}{假设}{假设}
\AtBeginDocument{
\renewcommand{\crefpairconjunction}{和}
\renewcommand{\crefmiddleconjunction}{、}
\renewcommand{\creflastconjunction}{和}
\renewcommand{\crefrangeconjunction}{至}
}
\floatname{algorithm}{算法}
% 中文过程名使用常规字形，避免切换到不支持中文的小型大写字形。
\renewcommand{\textproc}[1]{\textnormal{#1}}
\algrenewcommand{\algorithmicprocedure}{过程}
\algrenewcommand{\algorithmicif}{若}
\algrenewcommand{\algorithmicthen}{则}
\algrenewcommand{\algorithmicelse}{否则}
\algrenewcommand{\algorithmicend}{结束}
\algrenewcommand{\algorithmicreturn}{返回}

% 草稿辅助标记；投稿前删除所有用例。
\newcommand{\planned}[1]{{\color{gray}\itshape 【计划补充】#1}}
\newcommand{\todo}[1]{{\color{red}【待办：#1】}}

% !TEX root = ../main.tex
% 由分析脚本自动生成，请勿手动编辑。
% 论文中引用的所有数值均在此定义为宏。


% 沿用已确定的标题；仅在论文定位改变时重新讨论。
\title{\method： 以观测为锚点的跨轮次推测执行\\面向低延迟推理智能体}

\author{ACL 匿名投稿}

\begin{document}
\maketitle

% !TEX root = ../main.tex
\begin{abstract}
\planned{最后撰写。逻辑：推理智能体速度较慢，是因为每一轮都要等待上一轮的完整推理结束；然而，在上一轮推理仍在进行时，下一轮的输入往往已经确定。\method 基于目标模型自身最终会提交的动作及其在沙箱中获得的真实观测，提前启动下一轮，并且仅在下一轮输入完全一致时保留结果；报告 Qwen3-32B 和 Qwen3.5-9B 在 $\tau^2$-bench 上的核心结果。}
\end{abstract}


% !TEX root = ../main.tex
\section{引言}
\label{sec:intro}

\planned{按由现象逐步收束到机制、再展开到方法与证据的 V 形结构组织，以下各写一段。
(1) 现象：推理智能体完成一项任务需要数分钟；几乎所有时间都花在目标模型自身的生成上，且各轮次串行执行。
(2) 现有技术的局限：工具重叠推测执行~\citep{bai2026spork,nichols2025speculative} 的收益受工具耗时占比限制；跨步骤推测执行~\citep{ye2026speculative,xu2026leap} 必须在目标模型完成推理前预测其决策，因此需要针对工作负载训练的草拟器，或接受较低的一致率。
(3) 机制：目标模型在自身推理过程中就会显露决策；当历史不保留推理内容时，下一轮输入对当前轮的依赖仅来自可见输出和观测。循环等待的是确认，而非新信息。
(4) 方法：\method 及其三项策略。
(5) 证据与贡献：在实验部分完成后确定。}

% !TEX root = ../main.tex
\section{背景与问题设定}
\label{sec:background}

\subsection{推理智能体}
\label{sec:background:loop}

我们考虑一种由推理模型驱动、能够调用工具的智能体，并将该模型称为\emph{目标模型}（target）。
任务按轮次 $t=1,2,\dots$ 推进。
在第 $t$ 轮，目标模型读取输入 $\ctx{t}$，并生成 $Y_t=(\rsn{t},\vis{t},\act{t})$：其中，$\rsn{t}$ 是位于思考分隔符内的推理片段，$\vis{t}$ 是可选的可见文本，$\act{t}$ 是工具调用；若没有调用，则记为 $\act{t}=\bot$。
没有工具调用时，$\vis{t}$ 是一条回复，它结束当前智能体执行片段，并将控制权交还给用户。
工具调用在环境状态 $\env{t}$ 上执行：
\begin{equation}
  (\obs{t},\env{t+1}) = E(\env{t},\act{t}),
  \label{eq:env}
\end{equation}
随后，执行框架将本轮追加到已提交历史，并渲染下一轮输入：
\begin{equation}
  \hist_{t+1} = \hist_t \oplus \langle \vis{t},\act{t},\obs{t}\rangle,
  \qquad
  \ctx{t+1} = \render(\hist_{t+1}),
  \label{eq:history}
\end{equation}
其中，$\render$ 表示应用聊天模板并进行词元化。

\paragraph{历史策略。}
\Cref{eq:history} 丢弃已经完成的推理内容。
这遵循了模型开发者的建议：在多轮使用时，历史应仅保留模型每轮的最终输出，而不保留其思考内容~\citep{qwen2025qwen3card,qwen2026qwen35card}。
这一策略也与 \citet{bai2026spork} 的参考执行框架一致。
本文所有方法均采用这一策略。\footnote{对于最后一条用户消息之后的助手消息，如果调用方提供了推理内容，两种模型的参考聊天模板都会重新插入这些内容。我们的执行框架从不提供推理内容。“局限性”一节讨论了保留推理内容的执行框架。}
在此策略下，一轮执行仅通过其可见文本、动作及由此产生的观测影响后续执行。

\paragraph{延迟的来源。}
各轮次严格串行。
输入 $\ctx{t+1}$ 包含 $\obs{t}$，而 $\obs{t}$ 依赖于 $\act{t}$；后者只有在目标模型完成生成后才会确定。
因此，一轮的耗时等于完整生成时间 $R_t$ 加上工具延迟 $L_t$。
对于推理模型，$R_t$ 占据主导地位（\cref{sec:motivation}）。

\subsection{无损加速}
\label{sec:background:lossless}

参考执行即上述串行循环。

\begin{definition}[无损加速]
\label{def:lossless}
若一种加速执行在每项任务上提交的可见输出、动作和观测序列均与参考执行相同，并到达相同的最终环境状态，则称其为\emph{无损}（lossless）。
\end{definition}

这一定义约束完整轨迹。
它强于逐动作保证：逐动作保证只要求每个已提交动作均由目标模型生成，但执行轨迹仍可能发生分歧~\citep{xu2026leap}。
该定义还要求参考执行可复现。
因此，我们采用贪心解码，并使用批次不变内核（batch-invariant kernels）~\citep{he2025nondeterminism}，使给定输入的输出不依赖同一批次中的其他请求。

\subsection{自推测工具派发}
\label{sec:background:spork}

\spork~\citep{bai2026spork} 通过三个组件，将工具延迟隐藏在目标模型自身的推理过程中。
\textbf{(D1) 分叉探测。}目标模型解码第 $t$ 轮时，一个探测请求从相同前缀继续生成：前缀包含输入、已生成的推理，以及一个强制关闭推理并开始工具调用的后缀。
探测请求解码得到候选调用 $\hat a_t$。
由于复用了前缀缓存，其开销较低。
\textbf{(D2) 置信度门控。}只有当候选调用中工具名片段的最小词元概率不低于 $\theta$ 时，才会派发该候选。
否则，每再生成 $\Delta$ 个词元就重试一次探测，每轮最多探测 $K$ 次。
\textbf{(D3) 草稿复用。}当目标模型生成正式工具调用时，探测得到的词元被用作推测解码草稿~\citep{leviathan2023fast,chen2023accelerating}。

生成结束后，只有在 $\hat a_t\equiv\act{t}$ 时，才采用已派发调用的结果；这里的一致要求工具名和所有参数完全相同。
否则，正常执行 $\act{t}$。
我们使用默认参数的 \spork（$\theta=0.9$、$\Delta=50$、$K=5$）作为 \method 的底层机制及主要基线。
若一个正确候选在本轮开始后的 $d_t$ 时刻派发，\spork 最多节省 $\min(L_t,R_t-d_t)$。
因此，其收益受工具延迟限制；当工具很快时，这一收益较小。

% !TEX root = ../main.tex
\section{下一轮为何仍在等待}
\label{sec:motivation}

\planned{仅分析已有轨迹中的机制，不开展新的运行。
(a) 决策与审议之间的时间差：从首次得到高置信度且正确的探测结果，到推理结束之间的时间。
(b) 窄接口：在已启动、且动作匹配的后继轮次中，完整下一轮输入完全一致的比例。
(c) 时间花在哪里：每轮生成耗时与工具耗时，以及两种模型在锚点就绪后剩余时间的分布。
要点：串行循环等待的是确认，而非新信息。}

% !TEX root = ../main.tex
\section{\method}
\label{sec:method}

\Cref{sec:motivation} 表明，目标模型停止推理之前，下一轮输入往往早已确定。
\method 将这段剩余时间转化为实际进展。
一旦候选动作执行完毕、真实观测可用，\method 就启动完整的下一轮。
只有当最终确认的下一轮输入与参考执行原本会使用的输入完全一致时，才保留结果。
该方法包含启动阶段\stratLaunch，以及由\stratCommit 和\stratRecover 构成的提交阶段。
\Cref{fig:timeline} 将其时序与串行执行和工具重叠执行进行比较。

\begin{figure*}[t]
\centering
\colorlet{reasonfill}{blue!13}
\colorlet{toolfill}{orange!40}
\colorlet{nextfill}{blue!30}
\colorlet{succfill}{green!30}
\newcommand{\seg}[6]{%
  \draw[draw=black!55,fill=#5,rounded corners=1pt] ({#1},{#3-#4/2}) rectangle ({#2},{#3+#4/2});
  \node[font=\scriptsize] at ({(#1+#2)/2},{#3}) {#6};}
\begin{tikzpicture}[x=1.2cm,y=1cm]
  \node[anchor=east,font=\small] at (-0.1,0) {串行};
  \seg{0}{6}{0}{0.42}{reasonfill}{第 $t$ 轮：推理后调用（$R_t$）}
  \seg{6}{7}{0}{0.42}{toolfill}{工具（$L_t$）}
  \seg{7}{11}{0}{0.42}{nextfill}{第 $t{+}1$ 轮（$B_{t+1}$）}

  \node[anchor=east,font=\small] at (-0.1,-1.1) {工具重叠};
  \seg{0}{6}{-1.1}{0.42}{reasonfill}{第 $t$ 轮}
  \seg{1.5}{2.5}{-1.55}{0.3}{toolfill}{工具}
  \seg{6}{10}{-1.1}{0.42}{nextfill}{第 $t{+}1$ 轮}
  \node[font=\scriptsize,anchor=east] at (1.47,-1.55) {$d_t$};

  \node[anchor=east,font=\small] at (-0.1,-2.45) {\method};
  \seg{0}{6}{-2.45}{0.42}{reasonfill}{第 $t$ 轮}
  \seg{1.5}{2.5}{-2.9}{0.3}{toolfill}{工具}
  \seg{2.5}{6.5}{-3.3}{0.3}{succfill}{后继轮次 $t{+}1$（第二沙箱）}
  \node[font=\scriptsize,anchor=east] at (1.47,-2.9) {$d_t$};
  \node[font=\scriptsize,anchor=east] at (2.47,-3.3) {$\tau_t$};
  \draw[decorate,decoration={brace,amplitude=3pt,mirror}] (2.5,-3.52) -- (6,-3.52)
    node[midway,below=2pt,font=\scriptsize] {锚点就绪后的剩余时间 $\sigma_t$};
  \draw[very thick,black!70] (6.5,-2.2) -- (6.5,-3.5);
  \node[font=\scriptsize,anchor=west] at (6.55,-3.3) {$\contract_t=1$ 时提交};
  \draw[densely dotted,black!60] (10,-1.35) -- (10,-2.45);
  \draw[<->,>=Stealth,black!70] (6.55,-2.45) -- (10,-2.45)
    node[midway,above,font=\scriptsize] {节省：$\min(B_{t+1},\sigma_t)$};
\end{tikzpicture}
\caption{一个窗口的时序。工具重叠执行（\spork）在推理过程中执行高置信度预测的调用，从而隐藏工具延迟 $L_t$。\method 进一步在调用的真实观测于 $\tau_t$ 时刻就绪后立即启动下一轮，在第二沙箱中运行，并且仅在下一轮输入完全匹配时提交。相对于工具重叠执行，其节省时间为 $\min(B_{t+1},\sigma_t)$（\cref{prop:latency}）。图中时长仅作示意。}
\label{fig:timeline}
\end{figure*}

\subsection{锚点}
\label{sec:method:anchor}

\begin{definition}[锚点]
\label{def:anchor}
第 $t$ 轮的\emph{锚点}（anchor）是一个元组 $\anchor_t=(\hat a_t,\hat o_t,\hat S_{t+1},\tau_t)$，包含候选动作、执行该动作获得的观测与执行后状态，以及这些信息可用的时刻 $\tau_t$。
锚点必须满足以下三个性质。
\begin{enumerate}[label=(P\arabic*),leftmargin=*,nosep]
\item \textbf{因果性。}仅依据已提交历史 $\hist_t$、状态 $\env{t}$ 以及 $\tau_t$ 之前已经生成的 $Y_t$ 片段计算。
\item \textbf{隔离性。}$\hat a_t$ 在 $\env{t}$ 的私有副本上执行，因此计算锚点不会改变已提交状态。
\item \textbf{可验证性。}各组成部分均可与对应的已提交内容进行精确比较：动作通过规范化形式的相等性比较，观测通过摘要比较，执行后状态通过指纹比较。
\end{enumerate}
\end{definition}

启动阶段和提交阶段只读取这四个字段。
因此，任何满足 (P1)--(P3) 的来源都可以驱动这两个阶段，例如经过训练的动作草拟器~\citep{xu2026leap}，或存储历史轨迹的记忆~\citep{li2026specmemory}。
我们使用自推测派发来实例化锚点。
候选 $\hat a_t$ 就是 \spork 在 $d_t$ 时刻派发的调用。
\method 在 $\env{t}$ 的副本上执行该调用，而非将推测执行限制在只读工具上，因此预订或退款等改变状态的调用也可以作为锚点。
锚点在 $\tau_t=d_t+L_t$ 时刻就绪。
这一实例化方式既不需要训练，也不需要辅助模型。

\paragraph{下一轮输入的确定性。}
给定锚点后，可以在父轮次结束之前构造第 $t{+}1$ 轮的输入。
假设父轮次除工具调用外没有输出可见文本，\method 构造
\begin{equation}
  \hat X_{t+1} = \render\big(\hist_t \oplus \langle \varnothing, \hat a_t, \hat o_t\rangle\big).
  \label{eq:spec-input}
\end{equation}

\begin{proposition}[下一轮输入的确定性]
\label{prop:determinacy}
在\cref{eq:history} 的历史策略下：
(i) 若 $\vis{t}=\varnothing$、$\hat a_t\equiv\act{t}$ 且 $\hat o_t=\obs{t}$，则 $\hat X_{t+1}=\ctx{t+1}$；
(ii) 仅有 $\hat a_t\equiv\act{t}$ 和 $\hat o_t=\obs{t}$，并不能推出 $\hat X_{t+1}=\ctx{t+1}$。
\end{proposition}

第 (i) 点成立，因为\cref{eq:history} 使 $\ctx{t+1}$ 成为 $(\hist_t,\vis{t},\act{t},\obs{t})$ 的确定性函数。
对于第 (ii) 点，考虑工具调用相同、但可见前导文本不同的两个父轮次输出；它们会被渲染成不同输入。
此外，若历史保留 $\rsn{t}$，则 $\ctx{t+1}$ 将依赖在 $\tau_t$ 时刻尚未生成的推理内容。
因此，\cref{eq:history} 中的窄接口使提前构造输入成为可能，但仅凭动作相等不足以允许复用。
所以，\method 在提交时检查完整输入（\cref{sec:method:commit}）。

\subsection{启动阶段：\stratLaunch}
\label{sec:method:launch}

\paragraph{触发条件。}
当一个窗口的首个锚点就绪、且父轮次仍在推理时，\method 启动后继轮次；这里“仍在推理”是指输出既未关闭推理片段，也未开始工具调用。
若父轮次已到达这一位置，则剩余时间很少，窗口按底层机制继续执行。

\paragraph{第二沙箱。}
后继轮次以 $\hat X_{t+1}$ 为输入，在锚点执行后状态 $\hat S_{t+1}$ 的新副本 $\hat S'$ 中运行。
锚点的执行后状态必须保持不变，因为当 $\hat a_t\equiv\act{t}$ 时，底层机制会将其接纳为已提交状态。
因此，后继轮次的工具调用只修改 $\hat S'$。
即便是改变状态的调用，在通过验证之前也无法触及已提交环境。

\paragraph{后继轮次的执行。}
后继轮次由同一目标模型、使用相同生成配置执行，是一个完整轮次。
其内部仍启用自推测派发，探测预算为 $m$ 次（默认 $m=K$），因此后继轮次仍会重叠执行自身的工具调用。
我们考虑后继工具的两种策略。
默认的\emph{沙箱执行}（sandboxed）策略在 $\hat S'$ 中立即执行工具。
\emph{延迟执行}（deferred）策略则将工具执行推迟到提交决策之后，因此仅模型计算属于推测执行。
延迟执行适用于外部效果无法复制的工具，但可节省的时间更少。

\paragraph{窗口。}
各轮次被分成互不重叠的窗口 $(1,2),(3,4),\dots$。
后继轮次不会再启动自己的后继轮次，下一个窗口从已提交历史和状态开始。
窗口不会跨越用户消息或最终回复。
当轮次预算仅剩一轮时，仅运行底层机制。
深度为一的设计限制了浪费的工作量，并简化验证；更深或滑动窗口留待未来研究。

\subsection{提交阶段：\stratCommit 与\stratRecover}
\label{sec:method:commit}

如果父轮次以回复结束（$\act{t}=\bot$），则丢弃后继轮次。
否则，按照底层机制提交父轮次的调用，得到 $\obs{t}$、$\env{t+1}$ 和 $\ctx{t+1}$。
随后，\method 计算提交契约（commit contract）
\begin{equation}
\begin{aligned}
  \contract_t = \mathbf{1}\big[\,
    &\hat a_t\equiv\act{t} \;\wedge\; h(\hat o_t)=h(\obs{t}) \\
    &\wedge\; \phi(\hat S_{t+1})=\phi(\env{t+1}) \\
    &\wedge\; \hat X_{t+1}=\ctx{t+1} \;\wedge\; \hat g_{t+1}=g_{t+1}\,\big].
\end{aligned}
\label{eq:contract}
\end{equation}
动作检查沿用底层机制本身的规则。
摘要 $h$ 是观测文本的 SHA-256；输入相等实际上已蕴含该检查成立，但单独检查可给出精确的拒绝原因。
指纹 $\phi$ 对环境状态进行哈希，这里包括智能体数据库和用户数据库，从而确认后继轮次的起始状态与参考执行第 $t{+}1$ 轮的起始状态相同。
输入检查比较完整词元序列；\cref{prop:determinacy} 表明，该检查不能被动作相等替代。
最后，$g$ 汇总生成配置，包括模型、随机种子、词元预算、上下文窗口、思考标志、停止字符串和采样参数。

\paragraph{\stratCommit。}
若 $\contract_t=1$，\method 等待后继轮次完成，并再次计算 $\contract_t$。
随后，将后继轮次纳入历史，并将 $\hat S'$ 接纳为已提交状态。
后继轮次中与父轮次推理重叠执行的部分，即为节省的时间。

\paragraph{\stratRecover。}
若 $\contract_t=0$，或后继轮次失败，\method 就取消其生成请求，并等待正在执行的工具调用返回。
随后，在 $\ctx{t+1}$ 和 $\env{t+1}$ 上正常运行第 $t{+}1$ 轮。
排空（drain）操作确保参考路径恢复时，没有过期分支仍在运行或竞争 GPU。
因此，一次拒绝的代价包括排空时间，以及后继轮次与父轮次重叠期间造成的资源争用。
\Cref{alg:window} 概述了一个窗口的执行过程。

\begin{algorithm}[t]
\caption{\method 在第 $t$ 轮和第 $t{+}1$ 轮上的一个窗口。}
\label{alg:window}
\small
\begin{algorithmic}[1]
\Procedure{锚点就绪}{$\hat a_t,\hat o_t,\hat S_{t+1}$}
  \If{父轮次仍在推理}
    \State $\hat X_{t+1}\gets\render(\hist_t\oplus\langle\varnothing,\hat a_t,\hat o_t\rangle)$
    \State $\hat S'\gets\hat S_{t+1}$ 的副本 \Comment{第二沙箱}
    \State 在 $\hat S'$ 中以 $\hat X_{t+1}$ 启动后继轮次
  \EndIf
\EndProcedure
\Procedure{父轮次结束}{$Y_t$}
  \If{$\act{t}=\bot$}
    \State 取消并排空后继轮次；\Return
  \EndIf
  \State 按底层机制提交 $\act{t}$；得到 $\obs{t},\env{t+1},\ctx{t+1}$
  \If{存在后继轮次且 $\contract_t=1$}
    \State 等待其完成；重新检查 $\contract_t$
    \State 接纳其轮次及 $\hat S'$
  \Else
    \State 取消并排空后继轮次
    \State 在 $\env{t+1}$ 中以 $\ctx{t+1}$ 运行第 $t{+}1$ 轮
  \EndIf
\EndProcedure
\end{algorithmic}
\end{algorithm}

\subsection{保证}
\label{sec:method:guarantees}

下面说明 \method 在什么条件下满足\cref{def:lossless} 的无损性，以及能够节省多少时间。

\begin{assumption}[执行模型]
\label{ass:exec}
\leavevmode
\begin{enumerate}[label=(A\arabic*),leftmargin=*,nosep]
\item \textbf{确定性目标模型。}对固定输入和生成配置，目标模型提交到历史中的输出不依赖并发请求或注入的草稿。
\item \textbf{确定性且带指纹的环境。}$E$ 是状态和动作的确定性函数，且指纹相等蕴含状态相等。
\item \textbf{隔离与单次提交。}推测执行只接触私有副本；副本只能通过接纳操作成为已提交状态，且每轮至多接纳一次。
\item \textbf{原子提交。}在最终检查 $\contract_t$ 与接纳操作之间，已提交历史和状态不发生变化。
\end{enumerate}
\end{assumption}

\begin{proposition}[无损性]
\label{prop:lossless}
在 (A1)--(A4) 以及渲染局部性条件（\cref{app:proofs}）成立时，对任何满足 (P1)--(P3) 的锚点来源，\method 都是无损的。
\end{proposition}

\begin{proof}[证明概要]
证明对窗口进行归纳。
根据 (A1)，父轮次与参考轮次相同；根据 (A2) 以及底层机制的精确匹配规则，其提交的观测和状态也相同。
若 $\contract_t=1$，后继轮次与参考执行第 $t{+}1$ 轮具有相同的起始输入、配置和状态，因此由 (A1) 和 (A2)，它产生相同输出和相同工具状态转移。
接纳是其唯一生效方式 (A3)，且接纳时状态未发生变化 (A4)。
若 $\contract_t=0$，后继轮次没有留下任何已提交影响 (A3)，第 $t{+}1$ 轮完全按照参考执行运行。
\Cref{app:proofs} 给出完整证明。
\end{proof}

由此可得两点说明。
第一，\cref{eq:history} 的历史策略不是无损性的前提。
如果该策略不成立，则 $\hat X_{t+1}\neq\ctx{t+1}$，每个后继轮次都会被拒绝，\method 退化为底层机制；这一策略只决定 \method 能否获得收益。
第二，正确性不依赖锚点是否正确：错误锚点只会损失时间，不会损害保真性。
在随机解码下，不应对被拒绝的后继轮次重新采样，也不应按结果选择后继轮次，否则会改变已提交输出的分布（\cref{app:proofs}）。

\begin{proposition}[窗口延迟]
\label{prop:latency}
考虑锚点正确的窗口，设 $d_t\le R_t$，并忽略资源争用与管理开销。
令 $B_{t+1}$ 为底层机制下第 $t{+}1$ 轮的耗时，$\slack_t=(R_t-\tau_t)^+$ 为锚点就绪后的剩余时间。
相对于底层机制，\method 节省
\begin{equation}
  \Delta T_t=\min\big(B_{t+1},\slack_t\big).
  \label{eq:latency}
\end{equation}
相对于串行执行，底层机制节省 $\min(L_t,R_t-d_t)$，\method 节省 $\min(L_t+B_{t+1},R_t-d_t)$。
\end{proposition}

底层机制下，窗口耗时为 $\max(R_t,\tau_t)+B_{t+1}$。
使用 \method 时，耗时为 $\max(R_t,\tau_t+B_{t+1})$，两者之差即\cref{eq:latency}。
$\slack_t$ 是观测返回后仍在进行推理的时间；\citet{anon2026eta} 将其识别为推理智能体最大的可优化延迟预算。他们通过截断推理来缩短这段时间，而 \method 保留完整推理，并用下一轮的执行填充这段时间。
由此有两个推论。
第一，收益受限于锚点就绪后目标模型继续审议的时长，而非工具延迟，因此 \method 与工具重叠执行互为补充；工具越快，$\slack_t$ 甚至越大。
第二，窗口无法早于其中耗时最长的组成部分结束，因此相对于底层机制，每个窗口的加速比至多为 $2\times$。
资源争用和拒绝会降低实际收益。
\Cref{app:proofs} 分解这些成本，\cref{sec:analysis:regime} 将用实测结果检验\cref{eq:latency}。

% !TEX root = ../main.tex
\section{实验}
\label{sec:experiments}

\subsection{实验设置}
\label{sec:exp:setup}

\planned{$\tau^2$-bench 的航空与零售领域~\citep{barres2025tau2}；Qwen3-32B 和 Qwen3.5-9B~\citep{qwen2025qwen3,qwen2026qwen35}；启用批次不变内核的 vLLM~\citep{kwon2023vllm}；基线为串行执行与 \spork，后者采用与 \method 相同的候选隔离执行方式；明确披露历史策略。}

\subsection{RQ1：端到端延迟}
\label{sec:exp:latency}

\planned{逐任务配对延迟、P50 和 P95，以及任务级自助法（bootstrap）区间。}

\subsection{RQ2：无损性与任务质量}
\label{sec:exp:lossless}

\planned{轨迹完全一致性、最终状态相等性、奖励一致性，以及提交审计。}

\subsection{RQ3：机制}
\label{sec:exp:mechanism}

\planned{启动、接纳与拒绝统计；消融设置：延迟执行后继工具、后继轮次不探测、后继轮次仅探测一次，以及串行对照。}

% !TEX root = ../main.tex
\section{分析}
\label{sec:analysis}

\subsection{RQ4：\method 何时有效？}
\label{sec:analysis:regime}

\planned{验证延迟模型：比较逐任务预测节省时间与实际节省时间；分析不同模型（32B 与 9B）及模式（固定消息与交互）下的适用条件。}

\subsection{RQ5：成本}
\label{sec:analysis:cost}

\planned{额外生成的词元数、浪费的后继轮次工作量，以及沙箱开销。}

% !TEX root = ../main.tex
\section{相关工作}
\label{sec:related}

\planned{词元级推测解码~\citep{leviathan2023fast,chen2023accelerating}。工具重叠推测执行~\citep{nichols2025speculative,bai2026spork,li2026specmemory,chen2026aospec}。使用独立或经过训练的提议器进行跨步骤动作推测~\citep{ye2026speculative,xu2026leap,liu2026smc,zhong2026dualspec}。在环境副本中执行推测动作，是上述多个系统共有的做法，本文不将其视为贡献。能力对照表：提议器来源、是否需要训练、观测来源、验证粒度，以及轨迹级保证。}

% !TEX root = ../main.tex
\section{结论}
\label{sec:conclusion}

\planned{总结机制、方法和证据；关于更一般的观点，即推理智能体的延迟往往花在确认决策而非获取信息上，最多用一句话表述。}

% !TEX root = ../main.tex
\section*{局限性}

\planned{在跨工具调用保留推理内容的执行框架中，提交契约会拒绝所有后继轮次，因此应禁用 \method。每个窗口的前瞻深度为一。仅涉及两个领域和两个模型系列。每种配置仅运行一次。轨迹级无损性要求确定性解码。被拒绝的后继轮次带来额外计算成本。}


\bibliography{references}

\appendix
% !TEX root = ../main.tex
\section{证明与延迟模型}
\label{app:proofs}

\subsection{渲染局部性}
\label{app:proofs:locality}

提交时，\method 接纳后继轮次的历史，其中第 $t$ 轮的条目是推测条目 $\hat e_t=\langle\varnothing,\hat a_t,\hat o_t\rangle$，而非父轮次的实际条目 $e_t=\langle\vis{t},\act{t},\obs{t}\rangle$。
提交契约保证这两个条目渲染得到的第 $t{+}1$ 轮输入相同。
为了进一步推出它们在所有后续输入中也具有相同的渲染结果，我们采用以下条件。

\begin{enumerate}[label=(R),leftmargin=*,nosep]
\item \textbf{条目局部渲染。}历史渲染得到的文本，是各条目片段的拼接；每个片段由该条目自身内容以及历史中的角色序列确定。
\end{enumerate}

如果聊天模板依据每条消息的内容与角色渲染，并在消息边界放置特殊词元，则满足 (R)；两种模型的模板在移除推理内容后均如此。
相同词元序列解码得到相同文本，因此 $\hat X_{t+1}=\ctx{t+1}$ 蕴含 $\hat e_t$ 与 $e_t$ 生成相同片段。
由于两份历史的角色序列相同，(R) 进一步保证任意后续扩展都会得到相同输入。
如果实现选择在提交时重新追加父轮次的实际条目，就不需要 (R)。

\subsection{\texorpdfstring{\Cref{prop:lossless}}{无损性命题}的证明}
\label{app:proofs:lossless}

设参考执行提交条目 $e_1,e_2,\dots$，并经过状态 $\env{1},\env{2},\dots$。
如果两份历史的角色序列相同，且生成的片段相同，则称其\emph{等价}；由 (R)，等价历史在任意后续扩展下都会渲染为相同输入。
我们对窗口进行归纳，证明每个窗口开始时，\method 的已提交历史都与参考历史等价，且其已提交状态与参考状态相同。
初始情形成立，因为两种执行从相同消息和状态开始。

考虑从第 $t$ 轮开始的窗口。

\paragraph{父轮次。}
根据归纳假设，输入 $\ctx{t}$ 与参考输入相同。
由 (A1)，即使探测、后继轮次和注入草稿并发运行，输出 $Y_t$ 仍与参考输出相同。
若 $\act{t}=\bot$，则提交相同回复，并取消任何后继轮次；由 (A3)，后继轮次没有留下已提交影响，窗口结束时不变式仍成立。
否则，底层机制通过以下两种方式之一提交 $\act{t}$。
如果已派发候选满足 $\hat a_t\equiv\act{t}$，则接纳执行该候选的副本。
这一副本是在第 $t$ 轮期间从已提交状态复制的；由于本轮期间没有其他提交 (A3)，该状态等于 $\env{t}$。
环境对动作的依赖仅通过其规范化形式，因此由 (A2)，得到参考执行对应的 $(\obs{t},\env{t+1})$。
否则，直接在 $\env{t}$ 上执行 $\act{t}$，这就是参考状态转移。
在两种情况下，父轮次结束后的历史均与参考历史等价，状态均等于 $\env{t+1}$。

\paragraph{后继轮次被提交。}
假设已经启动后继轮次，且最终检查时 $\contract_t=1$。
于是 $\hat X_{t+1}=\ctx{t+1}$ 且 $\hat g_{t+1}=g_{t+1}$；同时，由 (A2)，$\phi(\hat S_{t+1})=\phi(\env{t+1})$ 推出 $\hat S_{t+1}=\env{t+1}$。
后继沙箱 $\hat S'$ 是 $\hat S_{t+1}$ 的副本，且只有后继轮次能够修改它。
由 (A1)，后继轮次产生参考输出 $Y_{t+1}$。
后继轮次内部的底层机制与父轮次一样提交调用，只是此时相对于 $\hat S'$ 执行；因此由 (A2)，其观测和执行后状态为参考执行的 $(\obs{t+1},\env{t+2})$。
接纳时，已提交历史成为 $\hist_t\oplus\hat e_t\oplus e_{t+1}$；根据\cref{app:proofs:locality}，它与参考历史等价；已提交状态则成为 $\env{t+2}$。
由 (A3)，没有其他影响进入已提交状态；由 (A4)，接纳时的状态就是检查过的状态。

\paragraph{后继轮次被拒绝或失败。}
如果没有启动后继轮次，或 $\contract_t=0$，或后继轮次失败，其影响均被限制在 $\hat S'$ 及其自身副本内 (A3)，排空后这些副本被丢弃。
随后，在 $\ctx{t+1}$ 和 $\env{t+1}$ 上运行第 $t{+}1$ 轮，父轮次的论证可原样应用于这一轮。

在所有情况下，窗口提交的条目均与参考执行一致，并在参考状态下结束，归纳完成。\qed

\paragraph{交互任务。}
使用模拟用户时，用户是环境的一部分。
窗口不会跨越用户消息，因此只要模拟器对相同对话的响应可复现，上述论证就适用于每个智能体执行片段。

\subsection{\texorpdfstring{\Cref{prop:latency}}{窗口延迟命题}的证明}
\label{app:proofs:latency}

假设锚点正确、$d_t\le R_t$，父轮次和后继轮次的耗时与底层机制下相同，且没有管理开销。
底层机制下，观测在 $\tau_t=d_t+L_t$ 时刻就绪，父轮次在 $\max(R_t,\tau_t)$ 时刻结束，随后第 $t{+}1$ 轮耗时为 $B_{t+1}$：
\[
T_{\text{sub}}=\max(R_t,\tau_t)+B_{t+1}.
\]
若 $\tau_t<R_t$，\method 在 $\tau_t$ 时刻启动后继轮次；当两个轮次都结束时，窗口结束：
\[
T_{\text{TA}}=\max(R_t,\tau_t+B_{t+1}).
\]
于是 $T_{\text{sub}}-T_{\text{TA}}=R_t+B_{t+1}-\max(R_t,\tau_t+B_{t+1})=\min(B_{t+1},R_t-\tau_t)$。
若 $\tau_t\ge R_t$，则不会启动后继轮次，节省时间为 $0=\min(B_{t+1},\slack_t)$。
这就证明了\cref{eq:latency}。

串行执行耗时为 $R_t+L_t+B_{t+1}$。
底层机制节省 $R_t+L_t-\max(R_t,d_t+L_t)=\min(L_t,R_t-d_t)$。
当 $\tau_t<R_t$ 时，\method 节省 $R_t+L_t+B_{t+1}-\max(R_t,d_t+L_t+B_{t+1})=\min(L_t+B_{t+1},R_t-d_t)$。
当 $\tau_t\ge R_t$ 时，两者节省的时间均为 $R_t-d_t$，上述公式仍成立。

对于加速比上界，令 $M=\max(R_t,\tau_t,B_{t+1})$。
两种情况下均有 $T_{\text{TA}}\ge M$，而 $T_{\text{sub}}\le 2M$，因此 $T_{\text{sub}}/T_{\text{TA}}\le 2$。

\subsection{实际收益}
\label{app:proofs:gain}

实际窗口在三个方面偏离理想模型：后继轮次与父轮次竞争 GPU，管理操作需要时间，且部分后继轮次会被拒绝。
令 $A=\max(R_t,\tau_t)$ 和 $B=B_{t+1}$ 分别为底层机制下父轮次与下一轮的耗时，并用撇号表示在 \method 下测得的耗时。

\paragraph{成功提交的窗口。}
设后继轮次在 $r\le A'$ 时刻启动、耗时为 $B'$；等待前的验证耗时为 $v_1$，接纳和清理耗时为 $v_2$。
窗口耗时为 $T_{\text{hit}}=\max(A'+v_1,r+B')+v_2$。
令重叠时间 $O=\min(B',(A'-r)^+)$，开销 $e=\max(A'+v_1,r+B')-\max(A',r+B')+v_2\in[v_2,v_1+v_2]$，则收益为
\[
G_{\text{hit}}=A+B-T_{\text{hit}}=O-(A'-A)-(B'-B)-e.
\]

\paragraph{被拒绝的窗口。}
设排空耗时为 $C_{\text{d}}$，回退轮次耗时为 $B^{\text{f}}$，其他管理开销为 $V$。
于是 $T_{\text{miss}}=A'+C_{\text{d}}+B^{\text{f}}+V$，且
\[
G_{\text{miss}}=-(A'-A)-(B^{\text{f}}-B)-C_{\text{d}}-V.
\]
若父轮次以回复结束，则 $B=B^{\text{f}}=0$，但仍需承担后继轮次造成的资源争用和排空成本。

\paragraph{期望收益。}
使用提交指示变量 $I$，有
\[
\begin{aligned}
\mathbb{E}[G]={}&\mathbb{E}[I\,O]-\mathbb{E}\big[\Delta A+I\,\Delta B+(1-I)\,\Delta B^{\text{f}}\\
&+I\,e+(1-I)(C_{\text{d}}+V)\big],
\end{aligned}
\]
其中 $\Delta A=A'-A$、$\Delta B=B'-B$，且 $\Delta B^{\text{f}}=B^{\text{f}}-B$。
资源争用通过 $\Delta A$ 和 $\Delta B$ 体现；禁用后继轮次内部的草稿复用也会影响这些项。
第一项必须保留联合期望：提交事件与较长重叠时间存在相关性，因此一般有 $\mathbb{E}[I\,O]\neq\Pr(I=1)\,\mathbb{E}[O]$。
如果每个成功提交的窗口至少收益 $g_{\min}>0$，每个其他窗口至多损失 $c_{\max}\ge 0$，那么提交率 $p>c_{\max}/(g_{\min}+c_{\max})$ 足以保证净收益为正。

\subsection{工具延迟执行策略}
\label{app:proofs:deferred}

在延迟执行策略下，后继轮次的全部工具执行，包括其自身的推测派发，都等待提交决策。
若提交，工具在 $\hat S'=\env{t+1}$ 中运行，并由 (A2) 产生参考状态转移，因此\cref{prop:lossless} 仍然成立。
若拒绝，则没有执行过任何后继工具。
由于工具延迟要在提交后支付，节省时间缩减为近似 $\min(G_{t+1},\slack_t)$，其中 $G_{t+1}$ 为后继轮次的生成时间。

\subsection{随机解码}
\label{app:proofs:stochastic}

采用采样时，将 (A1) 替换为以下要求：输出是输入、配置和请求随机种子的确定性函数；同时，后继轮次使用参考执行第 $t{+}1$ 轮原本会使用的种子。
归纳证明即可保持不变。
必须避免的是依赖结果的选择。
对被拒绝的后继轮次重新采样，或按内容、完成时间从多个后继轮次中进行选择，都会改变已提交输出的分布。
\method 不进行上述操作：被拒绝的后继轮次会被丢弃，第 $t{+}1$ 轮只在参考路径上生成一次。
我们的所有实验均采用贪心解码。

% !TEX root = ../main.tex
\section{实现细节}
\label{app:impl}

\planned{配置表；各模型的聊天模板与历史处理方式；Qwen3.5-9B 使用的参考循环内核；拒绝原因分类；沙箱分叉与接纳。}

% !TEX root = ../main.tex
\section{补充结果}
\label{app:results}

\planned{按领域和数据划分列出的结果表、逐任务延迟分布，以及交互模式结果。}


\end{document}
